\renewcommand{\Anno}{2006}

\renewcommand{\NumeroProblema}{2}

\renewcommand{\Testo}{Al Casin\`{o} di Montecarlo (comune della provincia di
	Lucca) va di moda il ``Gioco dei Pacchi''.  Un giocatore paga 35
	euro per una partita.  Di fronte al giocatore vengono posti 4
	pacchi chiusi, contenenti rispettivamente 10, 20, 30, 40 euro.
	Ovviamente il giocatore non sa il contenuto di ciascun pacco.  A
	questo punto il giocatore pu\`{o} aprire un pacco a sua scelta e,
	una volta visto il contenuto, pu\`{o} accettarlo o rifiutarlo: se
	accetta incassa la cifra contenuta nel pacco stesso e la partita
	finisce, se rinuncia pu\`{o} aprire un secondo pacco a sua scelta.
	Visto il contenuto del secondo pacco, il giocatore pu\`{o} ancora
	una volta accettare o rifiutare.  Se accetta incassa la cifra
	contenuta nel secondo pacco e la partita finisce, se rifiuta apre
	un terzo pacco, e questa volta \`{e} obbligato ad accettarne il
	contenuto.
	
	Due amici, Marco e Tullio, si apprestano a trascorrere una serata
	al Casin\`{o}, durante la quale giocheranno pi\`{u} volte al Gioco
	dei Pacchi.  Discutendo sulle strategie da attuare, Tullio
	afferma: ``Potendo aprire 3 pacchi, prima o poi quello da 30 o
	quello da 40 euro lo trovo.  Appena esce uno di questi 2 pacchi io
	accetto''.  Marco invece afferma: ``Io accetto solo se mi arriva
	il pacco da 40 euro, se proprio non esce vorr\`{a} dire che mi
	tengo il contenuto del terzo pacco''.
	
	\begin{enumerate}
		\item  Determinare, alla lunga, che cosa possono aspettarsi i 
		2 amici dalle loro strategie. 
	
		\item  E voi, come giochereste?
	\end{enumerate}}

\renewcommand{\Risposta}{\begin{enumerate}
		\item  Sia Tullio sia Marco incassano in media 35 euro a
		partita. Considerando il costo della partita, le 2 strategie
		risultano equivalenti e permettono ai 2 amici di concludere la
		serata in pareggio.
	
		\item  Esiste una strategia che consente un incasso medio
		superiore a 35 euro a partita: occorre accettare il primo
		pacco se e solo se \`{e} quello da 40 euro e poi eventualmente
		accettare il secondo pacco se e solo se \`{e} quello da 30 o
		40 euro.
	\end{enumerate}}

\renewcommand{\Aiutino}{\begin{enumerate}
	\renewcommand{\labelenumi}{(\alph{enumi})} 
	\item  Calcolare, sia per Tullio sia per Marco, la probabilit\`{a}
	di terminare la partita con 10, 20, 30, 40 euro. Calcolare quindi 
	l'incasso medio a partita facendo la media pesata dei vari incassi
	possibili.

	\item  Provare a vedere cosa succede se si rifiuta al primo colpo 
	il pacco da 30 euro.
\end{enumerate}}

\renewcommand{\SoluzioneUfficiale}{\TitoloSoluzione\DomandaA Tullio
statisticamente terminer\`{a} met\`{a} delle sue partite accettando il
pacco da 30 euro e met\`{a} delle sue partite accettando il pacco da
40 euro. Di conseguenza incasser\`{a} in media 35 euro a partita, pari
al costo della partita stessa. Pertanto, alla lunga, Tullio rimarr\`{a}
in pareggio.

Passiamo ora alla strategia di Marco. La probabilit\`{a} di non
trovare il pacco da 40 euro in 3 estrazioni \`{e} 1/4. Di conseguenza 
statisticamente Marco terminer\`{a} 3/4 delle sue partite incassando
40 euro. Nel rimanente 1/4 delle partite Marco incasser\`{a} con
uguale probabilit\`{a} 10, 20 o 30 euro, dunque in media 20 euro. Di
conseguenza, l'incasso medio di Marco sar\`{a} di
$$40\cdot\frac{3}{4}+20\cdot\frac{1}{4}=\frac{140}{4}=35$$
euro a partita, e cio\`{e} esattamente uguale a quello di Tullio.
Pertanto anche Marco, pur attuando una strategia diversa da quella di 
Tullio, alla lunga rimarr\`{a} in pareggio.

\DomandaB Mostriamo che esiste una strategia che consente un incasso
medio superiore ai 35 euro: tale strategia consiste nell'accettare il
primo pacco se e solo se contiene 40 euro, e poi eventualmente
accettare il secondo pacco se e solo se contiene 30 o 40 euro.

Calcoliamo l'incasso medio con questa strategia. 
\begin{itemize}
	\item  In 1/4 dei casi si incassano subito 40 euro con il primo
	pacco.

	\item  In 1/2 dei casi il primo pacco contiene 10 o 20 euro. La
	partita si concluder\`{a} dunque con il secondo o il terzo pacco
	non appena esce il pacco da 30 o quello da 40. Poich\'{e} i 2
	finali sono equiprobabili in questi casi l'incasso medio \`{e} di 
	35 euro.

	\item  Nel rimanente 1/4 dei casi il primo pacco contiene 30 euro.
	Dopo averlo rifiutato si aprono 3 possibilit\`{a}.
	\begin{itemize}
		\item  In 1/3 dei casi il secondo pacco sar\`{a} quello da 40 
		euro, che si incassano all'instante.
	
		\item  In 1/3 dei casi il secondo pacco sar\`{a} quello da 10 
		euro, che si rifiuta: a questo punto l'incasso sar\`{a} il
		valore del terzo pacco, che contiene in media (40+20)/2=30
		euro.
	
		\item  In 1/3 dei casi il secondo pacco sar\`{a} quello da 20 
		euro, che si rifiuta: a questo punto l'incasso sar\`{a} il
		valore del terzo pacco, che contiene in media (40+10)/2=25
		euro.
	\end{itemize}
\end{itemize}

Mettendo tutto insieme l'incasso medio risulta quindi
$$\frac{1}{4}\cdot 40+\frac{1}{2}\cdot 35+\frac{1}{4}\left(
\frac{1}{3}\cdot 40+\frac{1}{3}\cdot 30+\frac{1}{3}\cdot 25\right)=
\frac{425}{12},$$
e cio\`{e} poco pi\`{u} di 35.4 euro, ma in ogni caso pi\`{u} del
costo della partita.

Dimostriamo ora che questa strategia non \`{e} ulteriormente
migliorabile.

Intanto \`{e} evidente che il primo pacco va sicuramente accettato se 
contiene 40 euro, e rifiutato se contiene 10 o 20 euro. L'unico dubbio
sorge dunque di fronte ad un primo pacco da 30 euro.

Se si decide di accettarlo, si ricade sostanzialmente nella strategia 
di Tullio. Infatti in tal caso il secondo pacco si accetta se \`{e} da
40 (ovviamente), si rifiuta se \`{e} da 10 o 20 (infatti in tal caso
rimangono ancora i 2 pacchi da 30 e 40), e anche con il 30 si accetta 
(infatti in tal caso i 2 pacchi rimasti sono 10-40 oppure 20-40, ed in
entrambi i casi l'incasso medio al terzo tentativo sarebbe minore od
uguale di 30 euro).

Se invece si decide di rifiutare un primo pacco da 30 euro, la strategia \`{e}
sostanzialmente obbligata: \`{e} infatti evidente che un secondo pacco
da 40 euro va sicuramente accettato, mentre il calcolo svolto
precedentemente mostra che rifiutare un secondo pacco da 10 o 20 euro
porta ad incassarne in media, rispettivamente, 30 e 25 al terzo
tentativo, dunque di pi\`{u}. 

}

\renewcommand{\SoluzioniAlternative}{\TitoloSoluzioneAlternativaA 
}

\renewcommand{\Osservazioni}{
}

\renewcommand{\Commenti}{
}